% Generated draft: requires mathematical and source review.
\documentclass[11pt]{article}
\usepackage[margin=25mm]{geometry}
\usepackage{fontspec}
\setmainfont{Latin Modern Roman}
\usepackage{amsmath,amssymb,mathtools,unicode-math}
\setmathfont{Latin Modern Math}
\usepackage{enumitem,needspace,xurl,hyperref}
\setlength{\parindent}{0pt}
\setlength{\parskip}{5pt}
\setlength{\emergencystretch}{3em}
\newcommand{\N}{\mathbb{N}}
\newcommand{\Z}{\mathbb{Z}}
\newcommand{\Q}{\mathbb{Q}}
\newcommand{\R}{\mathbb{R}}
\newcommand{\C}{\mathbb{C}}
\newcommand{\F}{\mathbb{F}}
\newcommand{\A}{\mathbb{A}}
\newcommand{\PP}{\mathbb{P}}
\newcommand{\E}{\mathbb{E}}
\newcommand{\eps}{\varepsilon}
\newcommand{\dd}{\,\mathrm{d}}
\begin{document}
\section*{UnsolvedMath ID: 30001194}
\textbf{Source problem code:} catalogue entry
\subsection{Definitions and mathematical statement}
The record contains a complexity question and related rank-variety questions. The conventional setting below makes the complexity question precise, but the field and rank-variety conventions are reconstructed mathematical context rather than explicit hypotheses supplied by the record.

Let \(p\) be a prime, let \(k\) be an algebraically closed field of characteristic \(p\), and let \(\mu\) be a partition of an integer \(m\geq 1\). Write \(\mathfrak S_m\) for the symmetric group and \(S^\mu\) for its Specht module over \(k\): the submodule of the permutation module on \(\mu\)-tabloids generated by the polytabloids
\[
e_t=\sum_{\sigma\in C_t}\operatorname{sgn}(\sigma)\{\sigma t\},
\]
where \(t\) ranges over tableaux of shape \(\mu\), \(C_t\) is the subgroup permuting entries within each column, and \(\{t\}\) denotes the tabloid obtained by forgetting the order within each row.

Let \(\operatorname{core}_p(\mu)\) denote the partition obtained by successively removing rim hooks of length \(p\). The \(p\)-weight is
\[
w_p(\mu)=\frac{m-\lvert\operatorname{core}_p(\mu)\rvert}{p}.
\]
Call \(\mu\) a \(p\times p\) partition if every distinct part is divisible by \(p\) and the multiplicity of every distinct part is divisible by \(p\). Equivalently, its Young diagram can be tiled by aligned \(p\times p\) squares.

For a finite-dimensional \(k\mathfrak S_m\)-module \(M\), its complexity \(c_{\mathfrak S_m}(M)\) is zero if \(M\) is projective; otherwise it is the least integer \(c\geq 1\) for which the terms \(P_j\) of a minimal projective resolution satisfy \(\dim_k P_j=O(j^{c-1})\) as \(j\to\infty\).

The converse question in the supplied original statement has the following conventional formulation: for every prime \(p\), every algebraically closed field \(k\) of characteristic \(p\), and every partition \(\mu\) that is not \(p\times p\), must
\[
c_{\mathfrak S_m}(S^\mu)=w_p(\mu)
\]
hold? The supplied dated assessment describes the \(p\times p\) direction as settled and the converse as unproved in general; this is source-reported status, not an independently verified mathematical conclusion.

For context on the related questions, fix an elementary abelian subgroup \(E=\langle g_1,\ldots,g_r\rangle\cong(C_p)^r\) of \(\mathfrak S_m\), with \(r\geq1\). One conventional rank variety of a \(kE\)-module \(M\) is
\[
V_E^{\mathrm{rank}}(M)=\{0\}\cup\left\{\alpha\in k^r\setminus\{0\}:M\text{ is not free over }k[t]/(t^p),\quad t\cdot v=\sum_{i=1}^r\alpha_i(g_i-1)v\right\}.
\]
Its projectivization consists of the lines represented by its nonzero points, inside \(\mathbb P^{r-1}_k\).

For the relevant partitions whose parts are all divisible by \(p\), the record asks for a characterization of their Specht modules through rank varieties, asks whether an integer denoted \(n\) in a rank-variety formula is always \(1\), and asks whether a stated upper bound for the degree of the projectivized rank variety can be improved or made sharp. These questions cannot be formulated completely from the supplied record: it omits the relevant subgroup \(E\), any further restrictions on the partitions, the rank-variety formula defining \(n\), and the stated degree bound. The conventional definition above does not identify those missing data. The background also mentions a proposed universal bound on the complexity drop, but supplies neither that bound nor its hypotheses; no numerical inequality is reconstructed here.
\subsection{Short English statement}
In the conventional modular representation setting, the main question asks whether every partition that is not \(p\times p\) has a Specht module whose complexity equals its \(p\)-weight. Related questions concern rank varieties for partitions with parts divisible by \(p\), an unspecified integer \(n\), and a sharp degree bound. Their defining formula and hypotheses are missing. The supplied partially solved label is dated curation requiring independent review.
\subsection{Sources}
\begin{itemize}
\item Supplied local UnsolvedMath registry record, permanent ID \(30001194\), title \emph{Rank Varieties of Specht Modules}, source problem code \texttt{catalogue entry}, in the Oberwolfach Reports Open Problems collection. The record gives proposed year 2009 and names no individual proposer.
\item Supplied primary source citation: \emph{Mini-Workshop: Support Varieties} (2009), Oberwolfach Reports, pp. 549--592, DOI \texttt{10.4171/owr/2009/10}; \url{https://doi.org/10.4171/owr/2009/10}. The record attributes the questions to workshop open-problem discussions; no specific page or speaker is identified.
\item Supplied literature source labeled ``Hemmer 2009'': \url{https://doi.org/10.1007/s10801-009-0169-0}.
\item Supplied literature source labeled ``Lim 2009'': \url{https://doi.org/10.1016/j.jalgebra.2009.01.016}.
\item Supplied literature source labeled ``original Oberwolfach problem'': \url{https://doi.org/10.4171/owr/2009/10}.
\item Literature assessment dated exactly \texttt{2026-08-21}: the supplied curation labels the problem partially solved, says the cited Specht-module results settle substantial structured families and the \(p\times p\) direction, and says the converse and a proposed universal bound on the complexity drop are not proved in general. This is supplied dated curation derived from the OWR record, not independently verified literature assessment. The record itself requests independent verification of time-sensitive status and cited claims.
\item The registry labels the statement \texttt{corrected\_verified} and reports that its content was checked against the official Oberwolfach report, with notation, context, and question form normalized without changing the source-level claim. This verification claim is supplied metadata; the present draft does not independently confirm it or claim to have consulted the cited sources.
\end{itemize}
\end{document}
